% Folder 79, batch 05 — archivists' pages 81–100.
% Pass: Opus 5.5 (claude-opus-5-5), 2026-10-03 — first pass, unchecked against the pages by a human.
%
% His own pagination runs 40 to 50 on the rectos (p. 81 = 40, p. 83 = 41,
% p. 85 = 42, p. 87 = 43, p. 89 = 44, p. 91 = 45, p. 92 = 46, p. 94 = 47,
% p. 96 = 48, p. 98 = 49, p. 100 = 50); the pages between continue the text.
\documentclass[11pt,a4paper]{article}
\input{../preamble/grothendieck}

\folder{79}
\batch{5}
\pages{81}{100}
\foldertitle{2-polyèdres réguliers triangulés, et modules libres de rangs 2 sur les anneaux $\Lambda=\mathbb{Z}/n\mathbb{Z}$ et $\Lambda=\mathbb{Z}$ : notes manuscrites (s.d.).}
\dating{[à partir de 1980]}
\watermark{Édition de démonstration}

\begin{document}

\page{81}\note{p. 40 de l'auteur. La page s'ouvre au milieu d'une phrase commencée avant ce lot.}
peut prendre \ill{} $= \lambda\cdot(\lambda)$ avec $\lambda\in\Lambda$ quelconque
(\uncertain{il faut évid\supplied{emmen}t} $\lambda\in\Lambda^{*}$). Mais sauf erreur,
ces trois types de matrices engendrent $\mathrm{Sl}(2,\Lambda)$.
\marginal{à vérifier — en regard des deux lignes « ces trois types … $\mathrm{Sl}(2,\Lambda)$ », marquées d'un trait vertical}
On trouve ainsi (sauf erreur) que $u$ est l'identité sur $\widetilde{S}_M$ tout entier,
\struck{vu que $\mathrm{Sl}(2,\Lambda)$ est transitif}
Comme $\mathrm{Gl}_\Lambda(M)$ est transitif sur l'ens\supplied{emble} des repères, on trouve ainsi :

Si $\Lambda^{*}$ engendre $\Lambda$ additivement, \struck{alors} p.~ex. $\Lambda$ anneau
½ local, alors tout automorphisme de la structure
$\widetilde{P}_M$ $\bigl(\widetilde{S}_M,\ \widetilde{F}_M,\ T\colon \Lambda^{*}\to\mathfrak{S}_{\widetilde{S}_M}\bigr)$
provient d'un automorphisme unique de $M$.
\note{Ce paragraphe est marqué d'un trait vertical à gauche. Sous $\widetilde{F}_M$, un signe d'inclusion renversé renvoie à $\mathfrak{P}(\mathfrak{P}_3(S))$ : $\widetilde{F}_M$ est une partie de $\mathfrak{P}_3(S)$.}

Notons que pour une permutation $u$ de $\widetilde{S}_M = M^{*}$, dire qu'elle commute
à l'antipodisme et commute à la structure triangulaire signifie que l'on a
\[
\tag{21}
\begin{cases}
\text{a)}\ u(-x) = -u(x) & x\in M^{*}\\
\text{b)}\ u(x+y) = u(x)+u(y) & x,y \in M^{*},\ x\wedge y\in(\det M)^{*}
\end{cases}
\]
\note{après « $x,y$ », un mot biffé illisible.}
et dire qu'elle commute de plus à l'action de $\Lambda^{*}$ signifie qu'on a
(\ill{} de (21~a)), la condition plus forte
\[
\tag{22}
u(\lambda x) = \lambda\, u(x) \qquad x\in M^{*},\ \lambda\in\Lambda^{*}.
\]

\page{82}
Pour montrer que $u$ est induit par un automorphisme (unique) de $M$, il est tentant
d'écrire $M$, en termes de \emph{l'action} de
$\bigl(\widetilde{S}_M, \widetilde{F}_M, T\colon \Lambda^{*}\to\mathfrak{S}_{\widetilde{S}_M}\bigr)$, comme
\[
\tag{23}
M \simeq \text{quotient du \struck{$\mathbb{Z}$} $\mathbb{Z}$-module libre } \mathbb{Z}^{(M^{*})}
\]
par le sous-groupe engendré par les $[x+y]-[x]-[y]$ (pour $x,y$ \struck{\ill{}} $\in M^{*}$
tels que $x\wedge y\in(\det M)^{*}$) et par les $[-x]+[x]$. Désignons par $R_M$ ce
sous-groupe, une suite d'homomorphismes canonique
\[
\tag{24}
0 \longrightarrow R_M \xrightarrow{\ i\ } \mathbb{Z}^{M^{*}} \xrightarrow{\ p\ } M \longrightarrow 0
\]
et il faudrait montrer qu'elle est exacte (le seul p\supplied{rob}l\supplied{ème} étant
$\operatorname{Ker} p = \operatorname{Im} i$ ; il est clair que $pi = 0$). (Idem dans des
questions analogues pour \ill{} pour des modules libres de rang quelconque \add{$\geqslant 2$}.)
\marginal{NB La structure « triangulaire » + « antipod. » peut s'interpréter (en termes
de l'existence et de l'unicité des triangles complétant une arête) par une loi de groupe
\ill{} $(x,y)\mapsto x+y$ \uncertain{partiellement définie} \ill{}, alors par
$x\mapsto -x$, dont $M$ \uncertain{pourrait s'interpréter} comme (24) le groupe
\uncertain{enveloppant} …}
\note{note écrite en diagonale dans la marge gauche, en face de (23)–(24) ; lecture fragmentaire.}
Ainsi, $u$ satisfaisant à (21) induit une \struck{application} additive $\hat u$ de $M$
qui l'induit, et la condition (22) indique que ce dernier commute à $\Lambda^{*}$, donc
à $\Lambda$ si $\Lambda^{*}$ engendre $\Lambda$ additivement… Il est clair a

\page{83}\note{p. 41 de l'auteur.}
priori que cette condition supplémentaire est inutile pour indiquer que $\hat u$ est
$\Lambda$-linéaire, quand $\Lambda$ est un anneau absolu ($\simeq\mathbb{Z}$ ou
$\mathbb{Z}/n\mathbb{Z}$), puisque tout endomorphisme additif d'un $\Lambda$-module est
alors automatiquement $\Lambda$-linéaire. Par contre on comprend bien pourquoi, quand
$\Lambda$ n'est pas un anneau absolu, il faut absolument mettre une structure
supplémentaire sur $\widetilde{P}_M$ \add{$= (\widetilde{S}_M, \widetilde{F}_M)$} (genre
op\supplied{ération} de $\Lambda^{*}$ dessus) — en effet tout automorphisme
\emph{additif} (non nécessairement $\Lambda$-linéaire) de $M$ induit un automorphisme
de $\widetilde{P}_M$ ! Mais cette remarque ne s'applique pas telle quelle, si on tient
compte de la structure suppl\supplied{émentaire}
\[
\widetilde{F}_M = \coprod_i \widetilde{F}_{M,i}
\]
— \struck{sauf si} \add{si par exemple} $\Lambda^{*}=\{\pm1\}$ (p.~ex. $\Lambda =$
\struck{\ill{}} \uncertain{anneau de polynômes} sur $\mathbb{F}_2$, $\mathbb{F}_3$ ou
$\mathbb{Z}$), auquel cas l'ens\supplied{emble} d'indices $\Pi$ lui-même est réduit à un
élément !

\page{84}
J'en arrive maintenant au cas B) et je suppose dorénavant, pour ce cas,
\[
\tag{25}
\boxed{\,2\cdot 1_\Lambda \neq 0\,}
\]
de façon que \add{l'involution} $-\mathrm{id}_M$ opère librement sur
$\widetilde{S}_M = M^{*}$, sur $\widetilde{A}_M$ et sur $\widetilde{F}_M$ — on note
\struck{que} \uncertain{en effet} que
\[
\tag{26}
\{x,y,z\}\in\widetilde{F}_M \Longrightarrow x,y,z,-x,-y,-z \text{ sont tous distincts}.
\]
On pose
\[
\tag{27}
\begin{cases}
S_M = \widetilde{S}_M/\{\pm\mathrm{id}_M\} = M^{*}/\{\pm1\}\\
F_M = \widetilde{F}_M/\{\pm\mathrm{id}_M\} = \text{\struck{\ill{}}}
\end{cases}
\]
et on a une bijection canonique
\[
\tag{28}
F_M \xrightarrow{\ \sim\ } \bigl\{ f\in\mathfrak{P}_3(S_M) \bigm| \exists\, \tilde f\in\widetilde{F}_M,\ \text{avec } f = p(\tilde f)\bigr\}
\]
où
\[
\tag{29}
p\colon \widetilde{S}_M \to S_M,\quad x\longmapsto \dot x \quad\text{est l'application canonique.}
\]
On vérifie \struck{en effet} que si $f\in\mathfrak{P}_3(S_M)$ satisfait (28), alors il y a
exactement deux $\tilde f$ tels que $f = p(\tilde f)$, formant une orbite sous
$\{\pm\mathrm{id}_M\}$.

\page{85}\note{p. 42 de l'auteur.}
Par la suite on identifie toujours $F_M$ à une partie de $\mathfrak{P}_3(S_M)$.

On trouve, par passage aux quotients à partir \struck{(30)} de
$\vec{\widetilde{\tau}}_F$ \struck{$\widetilde{F}_M$} et $\widetilde{\tau}_F$ (5) (6),
des applications canoniques
\[
\tag{30}
\underset{\text{faces orientées}}{\vec{F}_M} \xrightarrow{\ \vec{\tau}_F\ } \vec{\Pi}_M = (\det M)^{*}
\qquad \vec{\tau}_F(\check{f}) = -\vec{\tau}_F(f)
\]
\[
\tag{31}
F_M \xrightarrow{\ \tau_F\ } \Pi_M = (\det M)^{*}/\{\pm 1_\Lambda\},
\]
ce qui permet pour $i\in\Pi_M$, de définir les $F_{M,i}$, d'où
\[
\tag{32}
F_M = \coprod_{i\in\Pi_M} F_{M,i}.
\]
On pose
\[
\tag{33}
A_M = \bigl\{ a\in\mathfrak{P}_2(S_M) \bigm| \exists f\in S_M,\ f\supset a \bigr\}
\]
\note{sic : $f\in S_M$, pour $f\in F_M$.}
et on trouve que l'on a aussi
\[
\tag{34}
A_M = \bigl\{ a\in\mathfrak{P}_2(S_M) \bigm| \exists\, \tilde a\in\widetilde{A}_M,\ \text{avec } a = p(\tilde a)\bigr\}
\]
— mais ici le nombre des $\tilde a$ \add{$\in\widetilde{A}_M$} au-dessus de $a$ est égal à 4
— si $a=\{s,t\}$, et $s=\dot x$, $t=\dot y$, les quatre choix sont
$\{x,y\}$, $\{-x,-y\}$ \struck{\ill{}}, $\{x,-y\}$, $\{-x,y\}$

\page{86}
qui se groupent en deux paires \struck{d'arêtes} \struck{\ill{}} formées d'arêtes
antipodiques de $\widetilde{P}_M$. Ces deux paires correspondent aux deux faces de
\add{\ill{}}
\[
P_M = (S_M, F_M)
\]
qui contiennent $a$ :
\begin{quote}
(35) Pour toute $a\in A_M$, il existe exactement \emph{deux} faces $f\in F_M$ telles
que $f\supset a$.
\end{quote}
On a de plus, pour mémoire :
\begin{quote}
(36) Soit $e\in\vec{\Pi}$, $i=\{\pm e\}$, soit
$\vec{F}_{M,e}\subset\vec{F}_{M,i}$,
$\vec{F}_{M,e} = \bigl\{\vec f\in\vec{F}_{M,i} \bigm| \vec{\tau}(\vec f) = e\bigr\}$,
alors $F_{M,e}$ est une « orientation » de $P_{M,i}$, i.e. si
$\vec f,\vec f'\in\vec{F}_{M,e}$ sont \emph{adjacentes}, i.e.
$|\vec f|\cap|\vec f'| = a\in\mathfrak{P}_2(S_M)$, alors les orientations de $a$
induites par $\vec f$, $\vec f'$ sont opposées.
\end{quote}
\marginal{NB d'où, interprétant $i$ comme un ens\supplied{emble} à 2 éléments
($\subset\vec{\Pi}$), on a $\mathrm{Or}(P_{M,i}) \simeq i$ (\uncertain{ens.}).}
On récupère les $F_{M,i}$ à partir des $A_{M,i}$ (mais non $F_M$ à partir de $A_M$) par
\[
\tag{37}
F_{M,i} = \bigl\{ f\in\mathfrak{P}_3(S_M) \bigm| \forall a\in\mathfrak{P}_2(f),\ \text{on a } a\in A_{M,i}\bigr\}.
\]

\page{87}\note{p. 43 de l'auteur.}
\struck{Pour étudier l'étoile d'un}
La propriété (35) montre que les \struck{$P_{M,i}$}
\[
\tag{38}
P_{M,i} = (S_M,\ A_{M,i},\ F_{M,i})
\qquad \text{\add{+ relations d'incidence (l'inclusion)}}\quad (i\in\Pi)
\]
sont des structures de 2-polyèdres triangulés, \uncertain{plus} qu'il n'est pas clair que
l'étoile de chaque sommet soit connexe (= contour). Choisissons un
$e\in i = \{\pm e\}\subset\vec{\Pi} = \det(M)^{*}$, on trouve par le contour combinatoire
orienté correspondant à un sommet $s=\dot x$ est donné par
\[
\tag{39}
\vec{\mathrm{Et}}_e(s) \simeq \text{contour défini par l'ensemble affine } L_x\subset M
\text{ formé des } y \text{ tels que } x\wedge y = e,
\]
espace homogène sous des translations $V_x = \{u\in M \mid x\wedge u = 0\} = \Lambda x$,
avec comme ens\supplied{emble} des sommets $L_x$, et le successeur de $y\in L_x$ étant
$y+x$.
\marginal{NB \struck{Pour tout $s\in S_M$,} \struck{le choix de $e\in i$ définit sur
$\mathrm{Et}(s)$ une structure de torseur sous $V_s$ de $M$, $V_s=\Lambda.(x)$
\ill{}} … \uncertain{définie par} $\Lambda.(-x)$ … $y+1$ …
structure d'orientation … opposées si $e$ est remplacé par $-e$. Donc
$\mathrm{Et}(s)$ est un \uncertain{torseur} sous $\Lambda$ \ill{} (\uncertain{$\Lambda$ tordu par $i$}).}
\note{note en diagonale dans la marge gauche, en grande partie biffée ; lecture très fragmentaire.}
On trouve donc :
\begin{quote}
(40) Les étoiles des points de $P_{M,i}$ sont connexes ssi $\Lambda_0=\Lambda$, i.e.
$\Lambda$ est un anneau absolu, $\Lambda\simeq\mathbb{Z}$ ou
$\Lambda\simeq\mathbb{Z}/n\mathbb{Z}$. Dans \add{tous les cas,} \struck{ce cas,} les
ordres des \struck{\ill{}} \add{fragments} sont tous égaux à \uncertain{card} $\Lambda$.
\end{quote}

\page{88}
On a donc une vraie structure polyédrale pour $P_{M,i}$ ssi $\Lambda=\Lambda_0$.
Quant à $P_M$ tout entier, c'est une vraie structure polyédrale ssi de plus
$\Pi_M$ se réduit à un élément, i.e.
\[
\tag{41}
\begin{cases}
P_M = (S_M, A_M, F_M) \text{ est un 2-polyèdre (i.e. étoiles connexes) ssi on a}\\
\Lambda\simeq\mathbb{Z} \text{ ou } \Lambda\simeq\mathbb{Z}/n\mathbb{Z},\ \text{avec } n\in\{3,4,6\}
\end{cases}
\]
\note{dans l'ensemble $\{3,4,6\}$ un premier chiffre est surchargé.}
(NB le cas $n=2$ est exclu d'emblée ici !).

En tous cas, on considère sur $S_M$ l'ens\supplied{emble} des structures triangulées
$(F_{M,i})_{i\in I}$, \struck{de} chacune pour son compte $P_M$. Sous l'hypothèse
\struck{que} $\Lambda$, \uncertain{autre} que $2\cdot 1_\Lambda\neq 0$. On aura pour
cette structure une suite exacte
\[
\tag{42}
1\longrightarrow \underbrace{\mathrm{Aut}_0(P_M)}_{\substack{\text{automorphismes de } P_M\\ \text{fixant chaque } F_{M,i}}}
\longrightarrow \mathrm{Aut}(P_M) \longrightarrow
\underbrace{\mathfrak{S}_{\Pi_M}}_{\text{pas mis surjectif}}
\]
\note{lecture de la légende sous $\mathfrak{S}_{\Pi_M}$ douteuse : « pas mis surjectif » ou « \uncertain{presque} surjectif ».}
et un homomorphisme de suites exactes

(43)

\begin{tikzcd}[column sep=small, row sep=small]
1 \arrow[r] & \mathrm{Sl}_\Lambda^{\prime}(M)^{\pm} \arrow[r] \arrow[d] & \mathrm{Gl}_\Lambda(M)^{\prime} \arrow[r] \arrow[d] & \mathbb{Z}_\Lambda^{\prime} \arrow[r] \arrow[d] & 1 \\
1 \arrow[r] & \mathrm{Aut}_0(P_M) \arrow[r] & \mathrm{Aut}(P_M) \arrow[r] & \mathfrak{S}_{\Pi_M} &
\end{tikzcd}

\marginal{à côté de la flèche $\mathbb{Z}^{\prime}_\Lambda\to\mathfrak{S}_{\Pi_M}$ : op\supplied{ère} sur $\Pi_M$ qui est un torseur sous $\mathbb{Z}^{\prime}_\Lambda$}
où $\mathbb{Z}^{\prime}_\Lambda = \Lambda^{*}/\{\pm 1_\Lambda\}$ défini moyennant
\uncertain{l'iso} (13) (p.~36).
\struck{$\mathbb{Z}^{\prime}_\Lambda =$}

\page{89}\note{p. 44 de l'auteur.}
\begin{quote}
(44) Le groupe $\mathrm{Gl}^{\prime}_\Lambda(M)^{\pm}$ est simplement transitif sur l'ens\supplied{emble}
des repères $(s,t,u)$ de $P_{M,i}$ ($i\in\Pi$ fixé), et $\mathrm{Sl}^{\prime}_\Lambda(M)$ est
simplement transitif sur l'ens\supplied{emble} des repères \emph{directs} de $P_{M,e}$
($e\in\vec{\Pi}$ fixé) — enfin, $\mathrm{Gl}^{\prime}_\Lambda(M)^{\varepsilon}$ est simplement
transitif sur l'ens\supplied{emble} des repères de $P_M$.
\end{quote}

\emph{Théorème.} Supposons que $\Lambda=\Lambda_0$ i.e. que les $P_{M,i}$ \add{($i\in\Pi_M$)}
\uncertain{soient} des vrais polyèdres. Alors

a) Pour $e\in\vec{\Pi}$, $\mathrm{Sl}^{\prime}_\Lambda(M)\xrightarrow{\ \sim\ }\mathrm{Aut}^{+}(P_{M,e})$

b) Pour $i\in\Pi$, $\mathrm{Gl}^{\prime}_\Lambda(M)^{\pm}\xrightarrow{\ \sim\ }\mathrm{Aut}(P_{M,i})$

c) \struck{Pour} $\mathrm{Gl}^{\prime}_\Lambda(M)\xrightarrow{\ \sim\ }\mathrm{Aut}(P_M)$
\note{à c), le mot « Pour » est biffé ; le texte ne fixe rien.}
Ainsi, $P_{M,i}$ étant défini, à isomorphisme près, comme le 2-polyèdre formé par $S_M$ et
l'un des structures $F_{M,i}$ de $\mathfrak{P}_3(S_M)$, i.e. l'ens\supplied{emble} des
structures de 2-polyèdre combinatoire \emph{régulier} sur l'ensemble $S_M$, associées aux
diverses structures ½symplectiques sur $M$.
\marginal{NB La connaissance de $F_M\subset\mathfrak{P}_3(S_M)$ permet déjà de reconstituer
$A_M$ (33) et par là relations d'incidences naturelles, $(F_M, A_M)$ est un graphe strict,
ayant $F_M$ comme ens\supplied{emble} des sommets, $A_M$ comme « arêtes » \ill{} (les
« sommets » sont tous d'ordre 3) — ceci posé, la décomposition $F_M = \coprod F_{M,i}$ est
\uncertain{la décomposition} en composantes connexes, \uncertain{ce qui} \ill{} la
décomposition précédente.}
\note{note en diagonale dans la marge gauche ; lecture partielle.}

\emph{Scholie.} On voit donc que l'on a des \emph{2-polyèdres} combinatoires
\emph{réguliers} \uncertain{associés} à des \ill{} anneaux, pour $\Lambda\simeq\mathbb{Z}$
ou $\Lambda\simeq\mathbb{Z}/n\mathbb{Z}$ fixé
(\uncertain{soit} $\Lambda_n$) …

\page{90}
\note{La page dresse trois correspondances, en colonnes : à gauche une classe de modules, à droite le groupoïde correspondant.}
\[
\tag{45}
\begin{array}{l}
\text{Modules symplectiques libres}\\
\text{de rang 2 sur } \Lambda\text{, « modulo »}\\
\text{symétrie } \pm1\\
\text{(groupoïde quotient …)}
\end{array}
\xrightarrow{\ \approx\ }
\begin{array}{l}
\text{groupoïde formé des 2-polyèdres}\\
\text{\struck{triangulés} orientés réguliers,}\\
\text{\uncertain{plongés} à un type \struck{\ill{}}}\\
\text{bien déterminé (dépendant de } n\text{)}
\end{array}
\]
\[
(M, e)\ \text{mod}\ {-\mathrm{id}_M} \longmapsto (S_M,\ F_{M,e}\subset F_{M,\{e,-e\}})
\]
\[
\tag{46}
\begin{array}{l}
\text{Modules ½symplectiques libres}\\
\text{de rang 2 sur } \Lambda\text{, « modulo »}\\
\text{symétrie } \pm1
\end{array}
\xrightarrow{\ \approx\ }
\begin{array}{l}
\text{groupoïde formé des \struck{connexes}}\\
\text{2-polyèdres triangulés réguliers}\\
\text{(orientables mais non orientés),}\\
\text{de type déterminé}\\
\text{(\uncertain{dépendant} de } n\text{)}
\end{array}
\]
\[
(M, i)\ \text{\struck{mod}}\ {\pm\mathrm{id}_M} \longmapsto (S_M,\ F_{M,i})
\]
\[
\tag{47}
\begin{array}{l}
\text{Modules libres de rang 2}\\
\text{sur } \Lambda\text{, « modulo »}\\
\text{symétrie } \pm1
\end{array}
\xrightarrow{\ \approx\ }
\begin{array}{l}
\text{groupoïde formé des \add{structures} \struck{\ill{}}}\\
\text{multi-2-polyédrales triangulées (paquets}\\
\text{de structures 2-polyédrales triangulées), de}\\
\text{type d'isomorphie \uncertain{bien} déterminé}\\
\text{(dépendant de } n\text{)}
\end{array}
\]
\[
M\ \text{mod}\ {-\mathrm{id}_M} \longmapsto \bigl(S_M,\ (F_{M,i})_{i\in\Pi_M}\bigr)
\]
Il s'agirait alors

a) de caractériser axiomatiquement les structures ainsi obtenues ;

b) d'expliciter \uncertain{ces} structures, \uncertain{en} associant aux
\struck{\ill{}} \add{$\mathbb{F}_2$-}\uncertain{gerbes} des modules $M$ dont la structure provient

\page{91}\note{p. 45 de l'auteur.}
c) d'expliciter la construction des $M$ eux-mêmes \struck{dont} provenant la structure
combinatoire envisagée ;

d) d'expliciter les \uncertain{autres} structures qui \uncertain{se} trouvent à partir de la
donnée combinatoire, telles : la structure de \struck{\ill{}} torseur sous
$\mathbb{Z}^{\prime}_\Lambda = \Lambda^{*}/\pm1$ sur l'ens\supplied{emble} $\Pi$,
\add{qui relève $\Pi$} \struck{l'application} la \uncertain{$\Lambda^{\pm}$}-torseur $\vec{\Pi}$
avec $\vec{\Pi} = \coprod_{i\in\Pi}\mathrm{Or}(P_{M,i})$ — mais il faut définir l'action de
$\Lambda^{*}$ sur $\vec{\Pi}$), l'application $\vec{\tau}_F\colon \vec{F}_M\to\vec{\Pi}$,
l'opération de $\mathbb{Z}^{\prime}_\Lambda$ sur $S_M$ …

e) Expliciter comment l'ens\supplied{emble} des structures polyédrales $P_{M,i}$ se déduit de l'une
d'entre elles

f) Examiner, \add{pour $\Lambda$ quelconque,} les structures combinatoires supplémentaires sur
les $P_{M,i}$ ou $P_M$, indiquant $\Lambda$ (p.~ex. l'opération de
$\mathbb{Z}^{\prime}_\Lambda$ sur $S_M$) qui permettent de \uncertain{énoncer} un théorème
analogue au précédent, et des reconstitutions de $M$ à symétrie $-1$ près …

\page{92}\note{p. 46 de l'auteur.}
\textbf{VI} Je commencerai par la question e) de la page précédente, qui avait déjà été
laissée en suspens, au §~III.

\emph{Proposition.} Soit $\Lambda$ un anneau quelconque (\uncertain{$M$ $\Lambda$-module libre de
rang 2}), considérons l'action de $\mathrm{Sl}_\Lambda(M)$ sur $M^{*}=\widetilde{S}_M$, de
$\mathrm{Sl}_\Lambda(M)^{\prime}$ sur $M^{*}/\pm1 = S_M$. On cherche \struck{dans le groupe
$\mathfrak{S}_M$} le commutant des actions \add{dans $\widetilde{\mathfrak{S}}_M$ resp.
$\mathfrak{S}_M$} : on va prouver que c'est l'image de $\Lambda^{*}$ resp. de
$\mathbb{Z}^{\prime}_\Lambda=\Lambda^{*}/\pm1$. Comme $M^{*}$ et \uncertain{$S_M$} …

En fait \uncertain{la question} : $S_M$ est \uncertain{un} espace homogène \uncertain{sous} le
groupe d'opérateurs $G$ ($\mathrm{Sl}_\Lambda(M)$ resp. $\mathrm{Sl}_\Lambda(M)^{\prime}$), le
commutant $Z$ dans $\mathfrak{S}_S$ de $G$ est isomorphe à $N_G(H)/H$, où $H$ est le
stabilisateur d'un point $s$ de $S$, \uncertain{un} \uncertain{isomorphisme} \struck{homomorphisme}
\[
Z \xrightarrow{\ \sim\ } N_G(H)/H
\]
\uncertain{étant} défini par la relation entre $Z$ et $N_G(H)$ définie par la formule
\[
\tag{48}
u(s) = g\cdot s \qquad (u\in Z,\ g\in N_G(H)).
\]
Calculons donc le normalisateur de $H$, dans les deux cas, en prenant une base $(x,y)$ de
$M$, et en prenant $s=x$ \uncertain{ou} $s=\dot x$.

\page{93}
\[
\tag{49}
\text{A)}\quad
\begin{cases}
G = \mathrm{Sl}(2,\Lambda) = \Bigl\{\begin{pmatrix} a & b\\ c & d\end{pmatrix}\in M_2(\Lambda) \Bigm| ad-bc=1\Bigr\}\\[4pt]
H = \Bigl\{\begin{pmatrix} 1 & t\\ 0 & 1\end{pmatrix} \Bigm| t\in\Lambda\Bigr\}
\end{cases}
\]
Soit $g$
\[
g = \begin{pmatrix} a & b\\ c & d\end{pmatrix}\in G, \quad\text{d'où}\quad
g^{-1} = \begin{pmatrix} d & -b\\ -c & a\end{pmatrix}\ \text{\ill{}}
\]
\[
(*)\quad \operatorname{int}(g)\begin{pmatrix} 1 & t\\ 0 & 1\end{pmatrix}
= \underbrace{\begin{pmatrix} a & b\\ c & d\end{pmatrix}\begin{pmatrix} 1 & t\\ 0 & 1\end{pmatrix}}_{\begin{pmatrix} a & at+b\\ c & ct+d\end{pmatrix}}
\begin{pmatrix} d & -b\\ -c & a\end{pmatrix}
= \begin{pmatrix} 1-act & a^2 t\\ -c^2 t & 1+act\end{pmatrix}
\]
donc
\[
\operatorname{int}(g)\begin{pmatrix} 1 & t\\ 0 & 1\end{pmatrix}\in H
\iff
\begin{cases} -c^2 t = 0\\ \text{\struck{\ill{}}}\ 1-act = 1\end{cases}
\quad\text{d'où } (1-ac\lambda)(1+ac\lambda) = 1
\]
Donc
\[
N_G(H) = \Bigl\{ g=\begin{pmatrix} a & b\\ c & d\end{pmatrix}\in\mathrm{Sl}(2,\Lambda) \Bigm|
\begin{array}{l} -c^2 t = 0\\ \text{\struck{\ill{}}}\ ac\,t = 0\end{array} \Bigr\}
\quad \text{\struck{\ill{}}}
\]
Or \struck{faisant $\lambda=1$, on trouve $c^2=0$, (donc \ill{})} comme $ad-bc=1$ i.e. $ad=1+bc$,
\uncertain{on voit} \add{comme $c^2=0\Rightarrow$} que $ad$ est inversible,
\struck{$(ad)^{-1}=1-bc$ \ill{}}
\[
(ad)^{-1} = 1-bc \qquad a^{-1} = (1-bc)\,d
\]
et comme $ac=0$, on doit avoir $c=0$.
Donc
\[
\tag{50}
N_G(H) = \Bigl\{ g = \begin{pmatrix} \lambda & b\\ 0 & \lambda^{-1}\end{pmatrix} \Bigm| \lambda\in\Lambda^{*},\ b\in\Lambda \Bigr\}
= \Bigl\{ g\in\mathrm{Sl}(2,\Lambda) \Bigm| \operatorname{int}(g)\begin{pmatrix} 1 & 1\\ 0 & 1\end{pmatrix}\in H \Bigr\}
\]
Donc \uncertain{nous} \struck{\ill{}} \uncertain{trouvons} \add{le groupe des} permutations de $M^{*}$
commutant à $\mathrm{Sl}(2,\Lambda)$ transforme $x$ en $\lambda x$ ($\lambda\in\Lambda^{*}$),
donc c'est l'homothétie $\lambda\,\mathrm{id}_M$, qui a la propriété \uncertain{sur} $x$ et commute
à $\mathrm{Sl}(2,\Lambda)$.

\page{94}\note{p. 47 de l'auteur.}
\marginal{On suppose maintenant $2\cdot 1_\Lambda\neq 0$}
\[
\tag{51}
\text{B)}\quad
\begin{cases}
G^{\prime} = \mathrm{Sl}(2,\Lambda)^{\prime} = G/\pm1\\
H^{\prime} = \text{image dans } G^{\prime} \text{ du groupe } H \text{ de (49), par } p\colon G\to G^{\prime}\ \text{(can.)}
\end{cases}
\]
Soit $g^{\prime}$ un élément de $G^{\prime}$, provenant de
\[
g = \begin{pmatrix} a & b\\ c & d\end{pmatrix}, \qquad ad-bc=1
\]
et considérons \struck{les} \struck{\ill{}}
\[
\operatorname{int}(g^{\prime})\cdot \underbrace{p\begin{pmatrix} 1 & t\\ 0 & 1\end{pmatrix}}_{h^{\prime}=p(h)}
= p\,\operatorname{int}(g)\begin{pmatrix} 1 & t\\ 0 & 1\end{pmatrix}
= p\begin{pmatrix} 1-act & a^2 t\\ -c^2 t & 1+act\end{pmatrix}
\]
on a
\[
\tag{52}
\operatorname{int}(g^{\prime})(h^{\prime})\in H^{\prime} \iff
\begin{cases} -c^2 t = 0\\ 1-act\in\pm1 \end{cases}
\]
Donc
\[
g^{\prime}\in N_{G^{\prime}}(H^{\prime}) \iff
\begin{cases} c^2 = 0\\ \forall t\in\Lambda,\ \text{on a } act=0 \text{ ou } act=2 \end{cases}
\]
\marginal{NB l'image de $G$ est formée des $\begin{pmatrix} \pm1 & \cdot\\ 0 & \cdot\end{pmatrix}$}
\note{lecture du NB en diagonale dans la marge gauche douteuse.}
Considérons le cas $t=1$, on aura donc
\[
ac = 0 \quad\text{ou}\quad ac = 2.
\]
Dans le premier cas, on a donc \uncertain{puisque} ($c^2=0\Rightarrow a\in\Lambda^{*}$) $c=0$,
ce qui exclut d'ailleurs $ac=2$ (puisque $2_\Lambda\neq0$ par hypothèse). Dans le deuxième cas,
on doit avoir, en itérant, \uncertain{on trouve}
\[
4_\Lambda = 0
\]
et de plus, la relation (52) \struck{\ill{}} s'écrit
\[
1-2t\in\pm1 \quad \forall t\in\Lambda \quad\text{i.e.}
\]
\[
\tag{53}
\forall t\in\Lambda,\ \text{on a } 2t=0 \text{ ou } 2t=2
\]
ce qui implique que l'idéal $2\Lambda$ est réduit à deux

\page{95}
éléments $2_\Lambda$ et $0_\Lambda = 4_\Lambda$. Notons que comme $c = a^{-1}\cdot 2\neq 0$
donc $a^{-1}2 = 2$, $a^{-1}$ inv. et $2\neq0$, on a \uncertain{donc}
\[
c = 2.
\]
Donc on a deux cas

1) Si \struck{$\Lambda$ est une \ill{}} \struck{$4_\Lambda=0$, $2_\Lambda\neq0$, et $2\Lambda$ est}
l'idéal $2\Lambda$ n'a que deux éléments $0$ et $2$ (i.e. $\Lambda$ est une extension de
l'une $\mathbb{F}_2$-algèbre \add{$\Lambda_0$} par un idéal engendré par $2$, et ayant 2 comme
seul élément non nul) alors
\[
\tag{54}
N_{G^{\prime}}(H^{\prime}) = \Bigl\{ p\begin{pmatrix} a & b\\ c & d\end{pmatrix} \Bigm|
\begin{array}{l} \text{\struck{\ill{}}}\ ad-bc=1\\ c\in\{0,2\}\end{array} \Bigr\}
\]
\[
N^{0}_{G^{\prime}}(H^{\prime}) \overset{\text{déf}}{=} \Bigl\{ p\begin{pmatrix} \lambda & b\\ 0 & \lambda^{-1}\end{pmatrix} \Bigm| \lambda\in\Lambda^{*},\ b\in\Lambda \Bigr\}
\]
\marginal{Ce cas est celui où $\Lambda$ est une \uncertain{extension} \uncertain{nilpotente} de
l'extension de $\mathbb{F}_2$ par $\mathbb{F}_2$, \uncertain{p.~ex.} $\mathbb{Z}/4\mathbb{Z}$
\ill{} \ill{} des \ill{} \uncertain{$\Lambda_0\to\mathbb{F}_2$}. En fait, si
$\sigma = \dot p\begin{pmatrix} 1 & 0\\ 2 & 1\end{pmatrix}$ alors $N_{G^{\prime}}(H^{\prime})$ est
½direct produit de $N^{0}_{G^{\prime}}(H^{\prime})$ par $\{1,\sigma\}$. Le $\sigma$ \ill{}
\uncertain{antipodisme} \uncertain{(55)} \ill{}.}
\note{note en diagonale dans la marge gauche ; lecture fragmentaire.}
admet $N^{0}_{G^{\prime}}(H^{\prime})$ comme sous-groupe d'indice 2.

2) Dans tout autre cas, \struck{(p.~ex. si $4_\Lambda\neq0$)},
\[
N_{G^{\prime}}(H^{\prime}) = N^{0}_{G^{\prime}}(H^{\prime})
\]
Dans ce cas \add{2°}, on voit donc que le commutant de $\mathrm{Sl}_\Lambda(M)^{\prime}$ dans
$\mathfrak{S}_{S_M} = M^{*}/\pm1$ est $\mathbb{Z}^{\prime}_\Lambda$.

Dans le cas 1°) par contre, $\mathbb{Z}^{\prime}_\Lambda$ est d'indice 2 dans ce commutant, qui
est isom. au produit direct (\add{pas seulement ½direct !}) de $\mathbb{Z}^{\prime}_\Lambda$ par
le sous-groupe $\{1,\sigma\}$. Le cas typique est celui de \struck{$\sigma(\dot x) = (x+2y)^{\cdot}$}
$\Lambda=\mathbb{Z}/4\mathbb{Z}$, alors $\mathbb{Z}^{\prime}_\Lambda = \{1\}$, mais le
groupe des automorphismes de $S_M$ \uncertain{commutant} \uncertain{aux} \uncertain{mouvements} de
l'octaèdre \add{$P_M$} qui commutent à $\mathrm{Aut}^{+}(P_M)$, n'est pas réduit à $\{1\}$ : il est
engendré par l'antipodisme de l'octaèdre.

\page{96}\note{p. 48 de l'auteur.}
Je peux reprendre maintenant le fil des réflexions du §~III. Soit
\[
\Lambda = \mathbb{Z}/n\mathbb{Z} \qquad\text{avec } n\in\mathbb{N}^{*},\ n\neq 2,4
\]
et considérons un module libre de rang 2 $M$ sur $\Lambda$, considérons une des structures
$P_{i_0} = P_{M,i_0}$ sur $S_M = S$ ($i_0\in\Pi_M$) — alors le groupe des automorphismes
$G = G_M$ de $P_{i_0}$ (\uncertain{sous-groupe} de $\mathfrak{S}_S$) ne dépend pas du choix de
$i_0$, vu \ill{} d'ailleurs le sous-groupe $G^{+}$ des autom.\ directs.
Soit $Z$ son commutant \struck{\ill{}}
\[
\tag{56}
\begin{cases}
Z = \mathrm{Comm.}_{\mathfrak{S}_S}(G^{+})\\
G = \mathrm{Aut}(P_{i_0}) \quad \text{(indépendant du choix de } i_0\in\Pi)
\end{cases}
\]
On a alors
\[
\tag{57}
\begin{cases}
\forall s,t\in S,\ \lambda\in Z, \text{ on a :}\\
\{s,t\}\in A_{i_0} \iff \{\lambda s, \lambda^{-1}t\}\in A_{i_0} \Longrightarrow t\neq \lambda\cdot s
\end{cases}
\]
\note{dans (57), au-dessus de « on a », une surcharge biffée illisible.}
et de plus
\begin{quote}
(58) $\forall\lambda\in Z$, l'unique structure 2-polyédrale triangulée sur $S$, \struck{\ill{}}
l'ens\supplied{emble} d'arêtes $A_\lambda$ défini par
\[
\{s^{\prime},t^{\prime}\}\in A_\lambda \iff \{s^{\prime},\lambda^{-1}t^{\prime}\}\in A_{i_0}
\]
est \struck{\ill{}} une des $A_i$ ($i\in\Pi$), soit $A_{\lambda\cdot i_0}$
\end{quote}
et \uncertain{ensuite} …

\page{97}
\begin{quote}
(59) L'application $\lambda\mapsto\lambda\cdot i_0$, $Z\to\Pi$ est bijective ; plus généralement,
si on \uncertain{considère} \ill{} l'application $\lambda\mapsto\lambda\cdot i$, pour un choix
de $i\in\Pi$, on trouve que l'ens\supplied{emble} des structures 2-polyédrales $\Pi$ sur $S$ est
un torseur sous $Z$.
\end{quote}
Il faudrait maintenant expliciter l'iso
\[
\tag{60}
Z \simeq \mathbb{Z}^{\prime}_\Lambda \simeq \Lambda^{*}/\{\pm1_\Lambda\}
\]
en termes de la seule structure combinatoire d'une des $P_i$, soit $P_{i_0}$.

Regardons pour \uncertain{ceci}, pour $\lambda\in Z$ fixé et $s\in S$ fixé, l'application
\[
\tag{61}
\begin{cases}
\mathrm{Et}_{P_{i_0}}(s) \xrightarrow{\ \sim\ } \mathrm{Et}_{P_{\lambda\cdot i_0}}(s)
\qquad (\lambda\cdot i_0 = i_1)\\
t \longmapsto \lambda\cdot t
\end{cases}
\]
Notons que les deux \uncertain{ensembles} sont des torseurs sous les groupes
\[
\tag{62}
\begin{cases}
\text{\struck{\ill{}}}\ \Gamma_{i_0} = \Lambda\wedge_{\mu_2}\mu(i_0)
\quad\text{resp.}\quad \Gamma_{i_1} = \Lambda\wedge_{\mu_2}\mu(i_1)\\
\text{où } \mu(i) = \text{ens.\ des orientations de } P_i
\end{cases}
\]

\page{98}\note{p. 49 de l'auteur.}
Ceci posé, je dis que
\begin{quote}
(63) L'application (61), $\mathrm{Et}_{P_{i_0}}(s) \xrightarrow[\sim]{\ \lambda_s\ } \mathrm{Et}_{P_{i_1}}(s)$
définie par $\lambda$ est un \uncertain{isomorphisme} de torseurs, pour un isomorphisme des
groupes d'opérateurs (uniquement déterminé bien sûr)
\[
\lambda_{\Gamma,s}\colon \Gamma_{i_0} \xrightarrow{\ \sim\ } \Gamma_{i_1},
\]
\end{quote}
\marginal{à vérifier}
Quand on a choisi des orientations de $P_{i_0}$ et de $P_{i_1}$, i.e.\ des \uncertain{origines}
des $\mu(i_0)$ et $\mu(i_1)$, alors $\lambda_{\Gamma,s}$ s'explicite par un isomorphisme de groupes
\[
\Lambda \xrightarrow{\ \sim\ } \Lambda
\]
donc par un élément $\gamma$ de $\Lambda^{*}$, \uncertain{pris} \struck{\ill{}} \add{d'origines}
dans l'ens\supplied{emble} des \struck{\ill{}} \add{torseurs} $\mu(i_0)$, $\mu(i_1)$, $\gamma$ est
remplacé par $-\gamma$. Ainsi, on \uncertain{peut} interpréter $\lambda_{\Gamma,s}$ comme un élément
\[
\tag{64}
\lambda_{\Gamma,s}\in
\Bigl[\underbrace{\Lambda\wedge_{\mu_2}\bigl(\mu(i_0)\wedge\mu(i_1)\bigr)}_{\Lambda\text{-module inversible } \Gamma_{i_0 i_1}}\Bigr]^{*}
\]
où, je le rappelle, $i_1 = \lambda\cdot i_0$.

On a d'autre part \struck{un \uncertain{isomorphisme} canonique} \add{\uncertain{un} \ill{}}
canonique
\[
\tag{65}
\Gamma_{i_0 i_1}^{*}/\{\pm1\} \simeq \Lambda^{*}/\pm1
\]
$\xi \longmapsto \xi$\ill{}
\note{l'exposant de $\xi$ ressemble à un 4 ; lecture non tranchée. Même signe en (66), p.~99.}

\page{99}
d'où enfin une application
\[
\tag{66}
\begin{aligned}
Z &\longrightarrow \mathbb{Z}^{\prime}_\Lambda = \Lambda^{*}/\{\pm1_\Lambda\}\\
\lambda &\longmapsto (\lambda_{\Gamma,s})^{\text{?}}
\end{aligned}
\]
\note{même exposant douteux qu'en (65) (un 4 ?), rendu ici par un point d'interrogation.}
qui a priori dépend du choix de $i_0\in\Pi$, et de $s\in S$. Mais on trouve
\begin{quote}
(67) L'application précédente est indépendante des choix de $s\in S$ et de $i_0\in\Pi$, et c'est
un \emph{isomorphisme} de groupes (en fait, c'est l'isomorphisme inverse de l'isomorphisme obtenu
p.~47').
\end{quote}
\marginal{à vérifier}
Ainsi, j'ai explicité, en termes d'une quelconque des structures $P_{M,i}$ — ou de la structure
$F_M$ toute entière, au choix — \struck{\ill{}} \uncertain{le} \uncertain{module} de ces structures,
l'opération de $\mathbb{Z}^{\prime}_\Lambda$ sur $S$, et par là, l'opération de
$\mathbb{Z}^{\prime}_\Lambda$ sur $\Pi$, qui fait de $\Pi$ un $\mathbb{Z}^{\prime}_\Lambda$-torseur.
\note{l'insertion « ou de la structure $F_M$ toute entière, au choix » est de sa main, au-dessus de la ligne ; « \uncertain{le module} » lu sous réserve (« \ill{} »).}

Explicitons maintenant le $\Lambda^{*}$-torseur $\vec{\Pi}^{*}$, on a une bijection canonique
\[
\tag{68}
\vec{\Pi} \simeq \coprod_{i\in\Pi}\bigl[\mathrm{Or}(P_{M,i}) = \mu(i)\bigr]
\]
\note{sous (68), une ligne biffée illisible ; en marge droite, « ens.\ des deux orientations de $P_i$ » glose $\mu(i)$.}

\page{100}\note{p. 50 de l'auteur. Ce feuillet est inachevé ; le raisonnement continue au-delà du lot.}
\uncertain{l'interprétation} \uncertain{qui nous} donne \uncertain{aussitôt} l'application
canonique de degré 2
\[
\tag{69}
\begin{aligned}
\vec{\Pi} &\xrightarrow[\text{degré } 2]{} \Pi\\
\{i,e\} &\longmapsto i \qquad (e\in\mu(i))
\end{aligned}
\]
Il faut \struck{définir} \add{décrire combinatoirement} une opération de $\Lambda^{*}$ sur
$\vec{\Pi}$, de telle façon que (69) soit compatible avec l'application canonique
\[
\Lambda^{*} \longrightarrow \mathbb{Z}^{\prime}_\Lambda = \Lambda^{*}/\{\pm1_\Lambda\}
\]
Soit donc $(i, e\in\mu(i)) \in\vec{\Pi}$, $\lambda\in\Lambda^{*}$, $\dot\lambda$ son image dans
$\mathbb{Z}^{\prime}_\Lambda$, on veut définir
\[
\tag{70}
\begin{cases}
\lambda(i,e) = (i^{\prime},e^{\prime}) \qquad e^{\prime}\in\mathrm{Or}(P_{i^{\prime}}) = \mu(i^{\prime})\\
i^{\prime} = \dot\lambda\cdot i,\quad e^{\prime} = \text{?}
\end{cases}
\]
\struck{On considérons l'application (61) ($s\in S$ choisi) $\mathrm{Et}_{P_{i_0}}(s)\to\mathrm{Et}_{P_i}(s)$ ;
le choix de $e$ définit une orientation du polygone $\mathrm{Et}_{P_{i_0}}(s)$, qui devient ainsi un
torseur sous $\mathbb{Z}/n\mathbb{Z}$ ; \add{par le choix d'une orientation} \ill{} de
$\mathrm{Et}_{P_i}$}
\note{le passage précédent est encadré et barré de deux traits obliques.}
Choisissons $s\in S$, alors
\[
\lambda_{\Gamma,s}\in\Lambda\wedge_{\mu_2}\bigl(\mu(i)\wedge\mu(i^{\prime})\bigr)
\xrightarrow{\ \text{degré } 2\ }\mathbb{Z}^{\prime}_\Lambda
\]
est au-dessus de $\dot\lambda$ ; \uncertain{donc} \ill{} $e^{\prime}\in\mu(i^{\prime})$,

\end{document}
